% Capítulo textual do trabalho.

\chapter{Conclusão}
\label{cap:conclusao}

A conclusão é uma das partes mais relevantes de um Trabalho de Conclusão de Curso~(TCC), pois é nela que o autor sintetiza os principais achados da pesquisa e apresenta suas reflexões finais. O objetivo principal da conclusão é proporcionar uma visão clara e objetiva sobre os resultados obtidos, além de contextualizá-los dentro da problemática abordada no trabalho. Para isso, uma conclusão eficaz deve seguir algumas diretrizes, garantindo que todos os pontos cruciais sejam abordados.

Primeiramente, a conclusão deve retomar os objetivos do trabalho. Ao revisitar as perguntas de pesquisa e os objetivos estabelecidos na introdução, o autor consegue avaliar se foram alcançados ou não. Este processo de revisão não apenas reforça a relevância do tema, mas também permite que o leitor entenda como as informações apresentadas ao longo do texto se conectam com o que foi inicialmente proposto. Segundo \citeonline{gil}, a conclusão deve ``revisitar os objetivos traçados, avaliando o quanto se conseguiu avançar na pesquisa''~\cite{gil}.

Em segundo lugar, é fundamental que a conclusão contenha uma síntese dos principais resultados e achados da pesquisa. Essa síntese deve ser objetiva e direta, evitando a inclusão de novos dados ou informações não discutidas previamente no corpo do texto. A análise crítica dos resultados obtidos pode ser feita de maneira a apontar as implicações e contribuições que o trabalho traz para o campo de estudo em questão. Como afirma~\citeauthoronline{lakatos}, ``a conclusão é o momento de apresentar as implicações do estudo realizado, bem como suas contribuições para a área de conhecimento''~\cite{lakatos}.

Por fim, a conclusão pode incluir sugestões para futuras pesquisas. Este aspecto é importante porque pode abrir novas possibilidades de investigação, além de indicar áreas que precisam de mais atenção ou aprofundamento. A inclusão de sugestões é uma forma de reconhecer que o conhecimento é um processo contínuo e que existem lacunas a serem preenchidas. Como destaca \citeonline{lakatos}, ``sugerir novos caminhos de pesquisa é uma maneira de contribuir para o desenvolvimento da área e para a construção do conhecimento''~\cite{lakatos}.

Em suma, a conclusão deve retomar os objetivos, sintetizar os resultados e sugerir novas pesquisas. Essa estrutura não apenas ajuda o leitor a entender o trabalho de forma mais clara, mas também enfatiza a contribuição do autor para a área estudada.
